\documentclass[preview]{standalone}

\usepackage{amsmath}
\usepackage{amssymb}
\usepackage{stellar}
\usepackage{definitions}
\usepackage{bettelini}

\begin{document}

\title{Stellar}
\id{linearalgebra-error-correcting-codes}
\genpage

\section{Error correction}

\begin{snippet}{error-correcting-codes-motivation}
    Suppose Alice wants to send a message to Bob. During transmission, some
    digits or bits may be changed. Bob would like to do two things:
    \[
        \begin{array}{ll}
            1. & \text{detect that an error occurred},\\
            2. & \text{correct the received message when possible}.
        \end{array}
    \]
\end{snippet}

\begin{snippetdefinition}{error-detecting-correcting-code-definition}{Error-detecting and error-correcting codes}
    An \emph{error-detecting code} is an encoding rule that allows the
    receiver to decide, for some received messages, that an error occurred.
    An \emph{error-correcting code} is an encoding rule that also allows the
    receiver to reconstruct the original message for some prescribed error
    patterns.
\end{snippetdefinition}

\section{Check digits}

\begin{snippetdefinition}{check-digit-procedure-definition}{A check digit procedure}
    Let
    \[
        C=C_{16}C_{15}\cdots C_2C_1
    \]
    be a string of \(16\) decimal digits, where \(C_1\) is the rightmost
    digit. Define
    \[
        p\colon\{0,1,\ldots,9\}\fromto\{0,1,\ldots,9\}
    \]
    by
    \[
        \begin{array}{c|cccccccccc}
            n & 0 & 1 & 2 & 3 & 4 & 5 & 6 & 7 & 8 & 9\\
            \hline
            p(n) & 0 & 2 & 4 & 6 & 8 & 1 & 3 & 5 & 7 & 9.
        \end{array}
    \]
    Equivalently, \(p(0)=0\), \(p(9)=9\), and for \(1\leq k\leq8\),
    \(p(k)\) is the remainder of \(2k\) modulo \(9\).

    Define
    \[
        S=C_1+p(C_2)+C_3+p(C_4)+\cdots+C_{15}+p(C_{16}).
    \]
    The string \(C\) passes the check if
    \[
        S\equiv 0\pmod{10}.
    \]
\end{snippetdefinition}

\begin{snippetexample}{check-digit-procedure-example}{A failed check digit}
    Consider
    \[
        C=1234\,5678\,9012\,3456.
    \]
    The check digit procedure gives
    \[
        S=64.
    \]
    Since
    \[
        64\not\equiv0\pmod{10},
    \]
    the string is detected as invalid.
\end{snippetexample}

\begin{snippet}{check-digits-limitations}
    Check digits are useful for detecting simple transcription errors, but
    they do not generally identify which digit should be changed. To correct
    errors systematically, it is natural to encode binary messages using
    linear algebra over a finite \field.
\end{snippet}

\section{Binary messages}

\begin{snippet}{binary-messages-explanation}
    Digital messages are represented by strings of bits, namely strings whose
    entries are \(0\) or \(1\). Text, images, videos and numbers can all be
    encoded into binary strings.
\end{snippet}

\begin{snippetexample}{integer-binary-expansion-37-example}{Binary expansion of \(37\)}
    Repeated division by \(2\) gives
    \[
        \begin{array}{rcl}
            37 &=& 2\cdot18+1,\\
            18 &=& 2\cdot9+0,\\
            9 &=& 2\cdot4+1,\\
            4 &=& 2\cdot2+0,\\
            2 &=& 2\cdot1+0,\\
            1 &=& 2\cdot0+1.
        \end{array}
    \]
    Reading the remainders from bottom to top,
    \[
        (37)_{10}=(100101)_2.
    \]
    Indeed,
    \[
        37=1\cdot2^5+0\cdot2^4+0\cdot2^3+1\cdot2^2+0\cdot2+1.
    \]
\end{snippetexample}

\begin{snippetdefinition}{binary-field-f2-definition}{The field \(\mathbb{F}_2\)}
    The \emph{binary field} is the \set
    \[
        \fieldtwo=\{0,1\}
    \]
    with addition and multiplication modulo \(2\). Addition is determined by
    \[
        0+0=0,
        \qquad
        0+1=1,
        \qquad
        1+0=1,
        \qquad
        1+1=0,
    \]
    and multiplication is determined by
    \[
        0\cdot0=0,
        \qquad
        0\cdot1=0,
        \qquad
        1\cdot0=0,
        \qquad
        1\cdot1=1.
    \]
\end{snippetdefinition}

\begin{snippetproposition}{binary-field-f2-is-field}{\(\mathbb{F}_2\) is a field}
    The \set \(\fieldtwo=\{0,1\}\), with addition and multiplication modulo
    \(2\), is a \field.
\end{snippetproposition}

\begin{snippetproof}{binary-field-f2-is-field-proof}{binary-field-f2-is-field}{\(\mathbb{F}_2\) is a field}
    The operations are ordinary integer addition and multiplication, reduced
    modulo \(2\). The additive identity is \(0\), and the multiplicative
    identity is \(1\). Since
    \[
        1+1=0,
    \]
    the element \(1\) is its own additive inverse. The only nonzero element is
    \(1\), and it has multiplicative inverse \(1\). Associativity,
    commutativity and distributivity are inherited from the integers modulo
    \(2\).
\end{snippetproof}

\section{Matrices over the binary field}

\begin{snippetdefinition}{binary-field-matrix-definition}{Matrices over \(\mathbb{F}_2\)}
    A \matrix over \(\fieldtwo\) is a \matrix whose entries belong to
    \(\fieldtwo=\{0,1\}\). If
    \[
        A\in\matrices_{m\times n}(\fieldtwo),
        \qquad
        B\in\matrices_{n\times p}(\fieldtwo),
    \]
    their product is defined by the usual formula
    \[
        (AB)_{ij}=\sum_{k=1}^{n}a_{ik}b_{kj},
    \]
    with all sums and products computed in \(\fieldtwo\).
\end{snippetdefinition}

\begin{snippetexample}{binary-field-matrix-addition-example}{Matrix addition over \(\mathbb{F}_2\)}
    In \(\fieldtwo\), one has \(1+1=0\). Hence
    \[
        \begin{pmatrix}
            1 & 0 & 1\\
            0 & 1 & 0\\
            0 & 0 & 1
        \end{pmatrix}
        +
        \begin{pmatrix}
            1 & 1 & 0\\
            0 & 0 & 1\\
            1 & 0 & 1
        \end{pmatrix}
        =
        \begin{pmatrix}
            0 & 1 & 1\\
            0 & 1 & 1\\
            1 & 0 & 0
        \end{pmatrix}.
    \]
\end{snippetexample}

\begin{snippetexample}{binary-field-matrix-vector-product-example}{Matrix-vector product over \(\mathbb{F}_2\)}
    Let
    \[
        A=
        \begin{pmatrix}
            0 & 0 & 1 & 0\\
            0 & 0 & 0 & 1\\
            1 & 1 & 0 & 0
        \end{pmatrix}
        \in\matrices_{3\times4}(\fieldtwo)
    \]
    and
    \[
        v=
        \begin{pmatrix}
            1\\
            1\\
            0\\
            1
        \end{pmatrix}
        \in\fieldtwo^4.
    \]
    Then
    \[
        Av=
        \begin{pmatrix}
            0\\
            1\\
            0
        \end{pmatrix}.
    \]
    Indeed, if \(A^j\) is the \(j\)-th column of \(A\), then
    \[
        Av=1A^1+1A^2+0A^3+1A^4.
    \]
    Since \(A^1=A^2=(0,0,1)^t\) and \(A^4=(0,1,0)^t\), the sum over
    \(\fieldtwo\) is
    \[
        (0,0,1)^t+(0,0,1)^t+(0,1,0)^t=(0,1,0)^t.
    \]
\end{snippetexample}

\section{Redundancy}

\begin{snippet}{redundancy-error-correction-idea}
    The basic idea behind error correction is to add redundancy to the
    transmitted message. A naive method would be to repeat a bit several
    times and decode by majority vote. For example, if Alice sends a bit five
    times and Bob receives
    \[
        01000,
    \]
    he may decide that the intended bit was \(0\). This corrects some errors,
    but it is inefficient: only one bit of information is transmitted in five
    positions, so the rate is \(1/5\).
\end{snippet}

\section{The Hamming code}

\begin{snippetdefinition}{hamming-7-4-parity-check-matrix-definition}{Parity-check matrix of the Hamming code}
    The Hamming \((7,4)\) code over \(\fieldtwo\) uses the parity-check
    \matrix
    \[
        H=
        \begin{pmatrix}
            1 & 0 & 1 & 0 & 1 & 0 & 1\\
            0 & 1 & 1 & 0 & 0 & 1 & 1\\
            0 & 0 & 0 & 1 & 1 & 1 & 1
        \end{pmatrix}
        \in\matrices_{3\times7}(\fieldtwo).
    \]
    The \(i\)-th column of \(H\) is the binary representation of \(i\):
    \[
        \begin{array}{c|c}
            i & \text{\(i\)-th column of \(H\)}\\
            \hline
            1 & (1,0,0)^t\\
            2 & (0,1,0)^t\\
            3 & (1,1,0)^t\\
            4 & (0,0,1)^t\\
            5 & (1,0,1)^t\\
            6 & (0,1,1)^t\\
            7 & (1,1,1)^t.
        \end{array}
    \]
\end{snippetdefinition}

\begin{snippetdefinition}{hamming-7-4-codeword-definition}{Codewords of the Hamming code}
    Let \(H\in\matrices_{3\times7}(\fieldtwo)\) be the parity-check \matrix
    of the Hamming \((7,4)\) code. A \emph{codeword} is a vector
    \[
        c\in\fieldtwo^7
    \]
    such that
    \[
        Hc=0.
    \]
    Equivalently, the \set of codewords is
    \[
        \{c\in\fieldtwo^7\suchthat Hc=0\}
        =
        \ltkernel L_H,
    \]
    where \(L_H\colon\fieldtwo^7\fromto\fieldtwo^3\) is the associated
    \lineartransformation \(L_H(c)=Hc\).
\end{snippetdefinition}

\begin{snippetdefinition}{hamming-7-4-encoding-definition}{Encoding in the Hamming \((7,4)\) code}
    In the Hamming \((7,4)\) code, the information message
    \[
        (a,b,c,d)\in\fieldtwo^4
    \]
    is encoded as the codeword
    \[
        C=
        \begin{pmatrix}
            a+b+d\\
            a+c+d\\
            a\\
            b+c+d\\
            b\\
            c\\
            d
        \end{pmatrix}
        \in\fieldtwo^7.
    \]
\end{snippetdefinition}

\begin{snippetproof}{hamming-7-4-encoding-definition-proof}{hamming-7-4-encoding-definition}{Encoding in the Hamming \((7,4)\) code}
    Put
    \[
        C=
        \begin{pmatrix}
            x\\
            y\\
            a\\
            z\\
            b\\
            c\\
            d
        \end{pmatrix}.
    \]
    The bits \(a,b,c,d\) are information bits, and \(x,y,z\) are redundancy
    bits chosen so that \(HC=0\). Multiplying by \(H\) gives
    \[
        HC=
        \begin{pmatrix}
            x+a+b+d\\
            y+a+c+d\\
            z+b+c+d
        \end{pmatrix}.
    \]
    Thus \(HC=0\) is equivalent to
    \[
        x+a+b+d=0,
        \qquad
        y+a+c+d=0,
        \qquad
        z+b+c+d=0.
    \]
    In \(\fieldtwo\), subtraction is the same as addition. Hence
    \[
        x=a+b+d,
        \qquad
        y=a+c+d,
        \qquad
        z=b+c+d,
    \]
    giving the displayed encoding formula.
\end{snippetproof}

\begin{snippet}{hamming-7-4-rate}
    The Hamming \((7,4)\) code transmits \(4\) information bits using \(7\)
    total bits. Its rate is
    \[
        \frac{\text{number of information bits}}{\text{number of transmitted bits}}
        =
        \frac{4}{7}.
    \]
\end{snippet}

\section{Syndrome decoding}

\begin{snippetdefinition}{hamming-syndrome-definition}{Syndrome}
    Let \(H\in\matrices_{3\times7}(\fieldtwo)\) be the parity-check \matrix
    of the Hamming \((7,4)\) code. For a received vector
    \[
        r\in\fieldtwo^7,
    \]
    the \emph{syndrome} of \(r\) is
    \[
        s=Hr\in\fieldtwo^3.
    \]
\end{snippetdefinition}

\begin{snippetexample}{hamming-7-4-zero-syndrome-decoding-example}{Zero syndrome}
    Let \(H\in\matrices_{3\times7}(\fieldtwo)\) be the parity-check \matrix
    of the Hamming \((7,4)\) code. If Bob receives
    \[
        r=
        \begin{pmatrix}
            x\\
            y\\
            a\\
            z\\
            b\\
            c\\
            d
        \end{pmatrix}
        \in\fieldtwo^7
    \]
    and \(Hr=0\), then \(r\) is a codeword. Under the assumption that at most
    one error occurred, Bob detects no error and decodes by removing the
    redundancy positions \(1,2,4\). The decoded message is
    \[
        (a,b,c,d),
    \]
    read from positions \(3,5,6,7\).
\end{snippetexample}

\begin{snippetproposition}{hamming-7-4-one-error-correction}{Correction of one error}
    Let \(H\in\matrices_{3\times7}(\fieldtwo)\) be the parity-check \matrix
    of the Hamming \((7,4)\) code. If a received word
    \(r\in\fieldtwo^7\) differs from a codeword by at most one coordinate,
    then the syndrome
    \[
        s=Hr
    \]
    identifies and corrects the error:
    \[
        s=0
        \Longrightarrow
        \text{no error is detected},
    \]
    while if \(s\neq0\), then \(s\) is a column of \(H\), and the error is
    in the corresponding position.
\end{snippetproposition}

\begin{snippetproof}{hamming-7-4-one-error-correction-proof}{hamming-7-4-one-error-correction}{Correction of one error}
    If \(r=C\), where \(C\) is a codeword, then
    \[
        Hr=HC=0.
    \]
    Suppose instead that there is exactly one error in position \(i\). Then
    \[
        r=C+e_i,
    \]
    where \(e_i\) is the \(i\)-th canonical vector of \(\fieldtwo^7\). Since
    \(C\) is a codeword,
    \[
        Hr=H(C+e_i)=HC+He_i=He_i.
    \]
    The vector \(He_i\) is the \(i\)-th column of \(H\). The columns of \(H\)
    are all nonzero and distinct, so the syndrome determines \(i\)
    uniquely. Correcting the error means changing the bit in position \(i\).
\end{snippetproof}

\begin{snippetexample}{hamming-7-4-seventh-position-correction-example}{Correcting the seventh position}
    Let \(H\in\matrices_{3\times7}(\fieldtwo)\) be the parity-check \matrix
    of the Hamming \((7,4)\) code. Suppose Bob receives \(r\in\fieldtwo^7\)
    and computes
    \[
        Hr=
        \begin{pmatrix}
            1\\
            1\\
            1
        \end{pmatrix}.
    \]
    This is the seventh column of \(H\). Therefore Bob corrects the received
    word by changing the seventh entry. After correction, he removes positions
    \(1,2,4\) and reads the information bits in positions \(3,5,6,7\).
\end{snippetexample}

\section{Limits of the code}

\begin{snippetexample}{hamming-7-4-multiple-error-limitation-example}{Two errors}
    Let \(H\in\matrices_{3\times7}(\fieldtwo)\) be the parity-check \matrix
    of the Hamming \((7,4)\) code. Suppose the correct codeword is \(C\), but
    two errors occur in positions \(i\) and \(j\), with \(i\neq j\). Then the
    received word is
    \[
        r=C+e_i+e_j.
    \]
    Its syndrome is
    \[
        Hr=HC+He_i+He_j=He_i+He_j.
    \]
    Since computations are in \(\fieldtwo\), the sum of two columns of \(H\)
    can coincide with another column of \(H\). Bob may then interpret the
    syndrome as a single error in the wrong position.
\end{snippetexample}

\begin{snippetproposition}{hamming-7-4-multiple-error-probability}{Probability of at most one error}
    Suppose that each transmitted bit is changed independently with
    probability \(p\). In a block of \(7\) bits, the probability of no errors
    is
    \[
        (1-p)^7,
    \]
    the probability of exactly one error is
    \[
        7p(1-p)^6,
    \]
    and the probability of at most one error is
    \[
        (1-p)^7+7p(1-p)^6.
    \]
    Therefore the probability of two or more errors is
    \[
        1-\bigl((1-p)^7+7p(1-p)^6\bigr).
    \]
\end{snippetproposition}

\begin{snippetproof}{hamming-7-4-multiple-error-probability-proof}{hamming-7-4-multiple-error-probability}{Probability of at most one error}
    No errors occur exactly when all \(7\) bits are transmitted correctly,
    which has probability \((1-p)^7\). Exactly one error can occur in any one
    of the \(7\) positions. For a fixed position, the probability is
    \(p(1-p)^6\), so the total probability of exactly one error is
    \(7p(1-p)^6\). These two events are disjoint, so their probabilities add.
    The event of two or more errors is the complement of the event of at most
    one error.
\end{snippetproof}

\begin{snippetexample}{hamming-7-4-multiple-error-probability-example}{A one percent bit error rate}
    If each bit is changed independently with probability
    \[
        p=0.01,
    \]
    then
    \[
        (1-p)^7+7p(1-p)^6\simeq0.998.
    \]
    Thus the probability of two or more errors in one block is approximately
    \[
        1-0.998=0.002,
    \]
    namely about \(0.2\%\).
\end{snippetexample}

\end{document}
